\documentclass{article}
\usepackage[utf8]{inputenc}
\usepackage{amsmath}

\begin{document}

\section*{2025 MIT Science Bowl High School Invitational - Round 15}

\subsection*{TOSS-UP}
\textbf{1) MATH} Short Answer Tim picks a random 3-digit number and squares it. What is the probability that its unit digit is a composite number?
\\ANSWER: $3/5$

\subsection*{VISUAL BONUS}
\textbf{1) MATH} Short Answer Your next bonus is visual. Edward and Anthony are travelling from city A to city B on the octahedral planet of Timlandia. Edward can only walk on the edges of the octahedron, while Anthony can travel over the surface of the planet. If Edward and Anthony both take their respective shortest possible paths, what is the ratio of their path lengths, in that order?
\\ANSWER: 4:3 (ACCEPT: $4/3$)

\subsection*{TOSS-UP}
\textbf{2) CHEMISTRY} Multiple Choice A 0.1 molar solution containing which of the following anions could be used to qualitatively separate a solution containing barium (II) and lead (II) ions?
\\W) Acetate
\\X) Hydroxide
\\Y) Fluoride
\\Z) Chloride
\\ANSWER: X) HYDROXIDE

\subsection*{VISUAL BONUS}
\textbf{2) CHEMISTRY} Short Answer Your next bonus is visual. Shown in the image are three organic molecules containing carbon-carbon single bonds, which have been indicated. Answer the following two questions about these bonds: 1) By number, rank the three indicated carbon carbon bonds by increasing bond length; 2) What is the name of the effect involving the delocalization of electrons in sigma orbitals responsible for the shortened carbon carbon bonds in these molecules?
\\ANSWER: 1) 1, 2, 3; 2) HYPERCONJUGATION

\subsection*{TOSS-UP}
\textbf{3) ENERGY} Short Answer Researchers in the Shalek Lab are using ATAC-seq (read: attack-seek) to study how stress alters transcriptional cell states in the liver. Which enzyme used in ATAC-seq isolates open chromatin and inserts sequencing adaptors, and is derived from mobile genetic elements?
\\ANSWER: TRANSPOSASE (ACCEPT: TN5 TRANSPOSASE)

\subsection*{BONUS}
\textbf{3) ENERGY} Multiple Choice The Quantum Engineering Group developed a fast algorithm using quantum annealing. What problem does their algorithm solve?
\\W) Finding roots of polynomials
\\X) Estimating non-analytic functions
\\Y) Inverting large matrices
\\Z) Globally minimizing functions
\\ANSWER: Z) GLOBALLY MINIMIZING FUNCTIONS

\subsection*{TOSS-UP}
\textbf{4) BIOLOGY} Short Answer Quinolones (read: quinn-o-lones) and camptothecins (read: camp-toh-theh-kins) are antibiotics that inhibit a certain enzyme to cause bacterial chromosomes to become toroidally shaped or plectonemic (read: pleck-toe-neem-ic) after replication. What general class of enzymes in DNA synthesis do they inhibit?
\\ANSWER: TOPOISOMERASES (ACCEPT: GYRASE)

\subsection*{VISUAL BONUS}
\textbf{4) BIOLOGY} Short Answer Your next bonus is visual. You are studying the DNA replication fork in a cell using FRET (read: fret). A donor fluorophore is attached to DNA helicase, and acceptor fluorophores are attached on DNA polymerases. Answer the following two questions: 1) Which of the following diagrams represents the signal on the leading strand? 2) Which of the following diagrams represents the signal on the lagging strand?
\\ANSWER: 1) E; 2) B

\subsection*{TOSS-UP}
\textbf{5) EARTH AND SPACE} Short Answer What structural feature is created by the rotational extension of crustal blocks in a back-arc extensional setting, typically seen as low-angle normal faults?
\\ANSWER: DETACHMENT FAULT

\subsection*{VISUAL BONUS}
\textbf{5) EARTH AND SPACE} Short Answer Your next bonus is visual. Answer the following two questions regarding the mineral shown in the photo. 1) What non-fluorescent silicate mineral is this specimen? 2) What crystal system does this mineral belong to?
\\ANSWER: 1) BERYL; 2) HEXAGONAL (ACCEPT: 1) RED BERYL, BIXBITE, DO NOT ACCEPT: 1) AQUAMARINE, MORGANITE, HELIODOR, MAXIXIE, EMERALD)

\subsection*{TOSS-UP}
\textbf{6) PHYSICS} Multiple Choice Which of the following statements is true regarding an ideal double-slit diffraction pattern, assuming that the slits have infinitesimal width?
\\W) The distance between consecutive peaks is inversely proportional to the energy of each photon
\\X) The central peak is the widest peak in the pattern
\\Y) The distance between consecutive peaks is proportional to the distance between the slits
\\Z) The intensity of the peaks increases with wavelength
\\ANSWER: W) THE DISTANCE BETWEEN CONSECUTIVE PEAKS IS INVERSELY PROPORTIONAL TO THE ENERGY OF EACH PHOTON

\subsection*{VISUAL BONUS}
\textbf{6) PHYSICS} Short Answer Your next bonus is visual. Shown in this image is a loop of wire consisting of a capacitor with initial charge Q and capacitance C, an inductor of inductance L, a resistor of resistance R, and an open switch connected in parallel. The switch is then closed. Answer the following two questions about the resulting circuit in terms of Q, C, and L: 1) For what value of R is the system critically damped? 2) A long time after the switch is closed, how much total energy has been dissipated by the resistor?
\\ANSWER: 1) $2\sqrt{L/C}$ (READ: TWO TIMES THE SQUARE ROOT OF L OVER C), 2) $Q^{2}/2C$ (READ: Q SQUARED OVER TWO C)

\subsection*{TOSS-UP}
\textbf{7) MATH} Short Answer Four Science Bowl teams compete in a single elimination bracket such that all four teams have distinct skill levels, and the team of higher skill always wins a match. After finals, another match is played between the losers to determine third place. If seeding is random, what is the probability that the final ranking matches the true order of skill levels?
\\ANSWER: $2/3$

\subsection*{BONUS}
\textbf{7) MATH} Short Answer Consider a differentiable function of that passes through the origin. Identify all of the following three statements that must be true about its inverse, given that it exists: 1) f inverse is continuous; 2) f inverse is differentiable; 3) f inverse passes through the origin.
\\ANSWER: 1 AND 3

\subsection*{TOSS-UP}
\textbf{8) CHEMISTRY} Short Answer In the Suzuki coupling, the insertion of a ligand from the boronate complex into the palladium metal center is an example of what type of organometallic mechanistic step?
\\ANSWER: TRANSMETALLATION

\subsection*{VISUAL BONUS}
\textbf{8) CHEMISTRY} Multiple Choice Your next bonus is visual. Shown in the image is the NMR spectrum of an unknown aromatic molecule that contains 8 hydrogen atoms. The integrations of the peaks from left to right are 2, 3, and 3, and the rightmost peak is a singlet. Which of the following is the most likely identity of the aromatic molecule?
\\W) Toluene
\\X) Methoxybenzene
\\Y) Benzyl alcohol
\\Z) para-methylphenol
\\ANSWER: X) METHOXYBENZENE

\subsection*{TOSS-UP}
\textbf{9) BIOLOGY} Short Answer Order the following three disorders of the eye in terms of increasing distance of the region affected to the retina: 1) Cataract; 2) Macular degeneration; 3) Conjunctivitus.
\\ANSWER: 2, 1, 3

\subsection*{VISUAL BONUS}
\textbf{9) BIOLOGY} Short Answer Your next bonus is visual. Shown in the image is an image of the heart of a patient suffering from low blood oxygenation. Answer the following three questions: 1) This image is produced with what technique, which uses ultrasound? 2) Identify the structure labeled by the red arrow, which is a temporary hole between the two atria present at birth; 3) The failure of the previous hole to close causes which of the labeled chambers, A or B, to have less oxygenated blood than usual?
\\ANSWER: 1) ECHOCARDIOGRAPHY; 2) FORAMEN OVALE (ACCEPT: PATENT FORAMEN OVALE) 3) B.

\subsection*{TOSS-UP}
\textbf{10) PHYSICS} Short Answer The Ehrenfest theorem relates the expected values of the position and momentum operators for a quantized system. The Ehrenfest theorem is analagous to what law in classical mechanics?
\\ANSWER: NEWTON'S SECOND LAW (ACCEPT: DEFINITION OF MOMENTUM)

\subsection*{BONUS}
\textbf{10) PHYSICS} Short Answer Ethan and Katie are on two spherical planets of uniform density. The moment of inertia of Katie's planet is 24 times that of Ethan's planet, while the surface gravitational field of Katie's planet is 1.5 times that of Ethan's planet. To the nearest tenth, what is the ratio of the escape velocity of Katie's planet to Ethan's planet?
\\ANSWER: 1.7

\subsection*{TOSS-UP}
\textbf{11) ENERGY} Multiple Choice Researchers at CSAIL are studying Hebbian descent, a technique designed to overcome vanishing gradients in gradient descent. Which of the following explains why vanishing gradients are undesirable in gradient descent?
\\W) Vanishing gradients prevent networks from updating weights to learn
\\X) Vanishing gradients lead to floating point errors in calculating weights
\\Y) Vanishing gradients require additional iterations to achieve convergence
\\Z) Vanishing gradients cause networks to overfit weights to input data
\\ANSWER: W) VANISHING GRADIENTS PREVENT NETWORKS FROM UPDATING WEIGHTS TO LEARN

\subsection*{BONUS}
\textbf{11) ENERGY} Short Answer Scientists at the Nelson Lab used terahertz light to create a metastable magnetized state. Identify all of the following three materials that are also considered metastable: 1) Supercooled water; 2) Technetium-99m; 3) Rydberg atom.
\\ANSWER: ALL

\subsection*{TOSS-UP}
\textbf{12) EARTH AND SPACE} Multiple Choice The first stage of mid-latitude cyclone formation involves a low-pressure center emerging along what type of front?
\\W) Warm
\\X) Cold
\\Y) Stationary
\\Z) Occluded
\\ANSWER: Y) STATIONARY

\subsection*{BONUS}
\textbf{12) EARTH AND SPACE} Short Answer Identify all of the following three statements that are true of plastic flow in glaciers: 1) Plastic flow dominates cold-based glaciers; 2) Plastic flow operates effectively throughout the entire glacier; 3) Plastic flow requires glacial ice to be at or above the pressure melting point.
\\ANSWER: 1 ONLY

\subsection*{TOSS-UP}
\textbf{13) MATH} Short Answer A circle in the cartesian plane is centered at (1,8). If the circle is tangent to a line at (0, 6), what is the x coordinate of the x intercept of the line?
\\ANSWER: 12

\subsection*{VISUAL BONUS}
\textbf{13) MATH} Short Answer Your next bonus is visual. In graph theory, an Eulerian path is a path that visits every edge of a graph exactly once. The path may start and end on different vertices, and revisiting vertices is allowed. Identify, by number, all of the displayed graphs that have an Eulerian path.
\\ANSWER: 1 AND 3

\subsection*{TOSS-UP}
\textbf{14) ENERGY} Multiple Choice Chemists at the Drennan Lab used data processing to artificially enhance the structure of carbon monoxide dehydrogenase. Which of the following experimental methods did they use to determine the initial structure of carbon monoxide dehydrogenase without crystallizing the protein sample?
\\W) X-ray diffraction
\\X) Cryogenic electron microscopy
\\Y) Scanning electron microscopy
\\Z) Nuclear magnetic resonance
\\ANSWER: X) CRYOGENIC ELECTRON MICROSCOPY

\subsection*{VISUAL BONUS}
\textbf{14) ENERGY} Short Answer Your next bonus is visual. The image shows free-range atoms in ultracold quantum gases studied by the Zwierlein Group. By letter, identify all of the following four particles that could be on the left side of the image: a) Free electron; b) Helium-3 atom; c) Helium-4 atom; d) Sodium-23 atom.
\\ANSWER: C AND D (ACCEPT: HELIUM-4 AND SODIUM-23)

\subsection*{TOSS-UP}
\textbf{15) EARTH AND SPACE} Short Answer Order the following four elements by increasing mass fraction of the observable universe: 1) Oxygen; 2) Carbon; 3) Lithium; 4) Iron.
\\ANSWER: 3, 4, 2, 1

\subsection*{BONUS}
\textbf{15) EARTH AND SPACE} Short Answer Order the following three fundamental forces in chronological order of when they first became independent forces in the inflationary universe hypothesis: 1) Weak force; 2) Strong force; 3) Gravitational Force
\\ANSWER: 3, 2, 1

\subsection*{TOSS-UP}
\textbf{16) PHYSICS} Multiple Choice In the 1967 novel One Hundred Years of Solitude, José Arcadio Buendía attempts to use a permanent iron magnet to collect gold along a riverbank. Which of the following explains why the magnet failed?
\\W) Gold has a high Curie temperature
\\X) Gold is very paramagnetic
\\Y) Gold has a low magnetic permeability
\\Z) Gold has a low magnetic susceptibility
\\ANSWER: Z) GOLD HAS A LOW MAGNETIC SUSCEPTIBILITY

\subsection*{BONUS}
\textbf{16) PHYSICS} Multiple Choice For a free particle traveling through a constant potential, which of the following pairs of observables can be simultaneously measured?
\\W) Position and momentum
\\X) Position and angular momentum
\\Y) Energy and position
\\Z) Energy and momentum
\\ANSWER: Z) ENERGY AND MOMENTUM

\subsection*{TOSS-UP}
\textbf{17) BIOLOGY} Short Answer Fluoroacetic acid (read: floor-o-uh-cetic) is highly toxic because it interferes with a particular enzyme in the Krebs cycle, resulting in the buildup of citrate. What enzyme does fluoroacetic acid inhibit?
\\ANSWER: ACONITASE

\subsection*{BONUS}
\textbf{17) BIOLOGY} Short Answer Alice's fruit fly embryo sample and Bob's basidiomycete (read: buh-SID-ee-oh-MY-seat) sample both contain cells that possess multiple nuclei. Identify all of the following three statements that are true about their samples: 1) The formation of both samples could be recreated by applying a microtubule depolymerization inhibitor to a cell; 2) Bob's samples are collectible through the majority of the basidiomycete lifetime; 3) If Alice and Bob each sequence the DNA within the nuclei of their own samples, Alice would observe less SNPs within her sample compared to Bob.
\\ANSWER: 2 AND 3

\subsection*{TOSS-UP}
\textbf{18) CHEMISTRY} Multiple Choice Which of the following compounds are most likely to be directly produced by the gas-phase reactions between alkenes and ozone in Earth's atmosphere?
\\W) Alkanes
\\X) Alkynes
\\Y) Aldehydes
\\Z) Ethers
\\ANSWER: Y) ALDEHYDES

\subsection*{BONUS}
\textbf{18) CHEMISTRY} Short Answer Identify all of the following three molecules that exhibit 3-fold symmetry: 1) Methane; 2) Phosphorus pentachloride; 3) Sulfur hexafluoride.
\\ANSWER: ALL

\subsection*{TOSS-UP}
\textbf{19) MATH} Multiple Choice A list of numbers X contains one 1, two 2s, and so on, where the number n appears n times for all n from 1 to 2025. Which of the following statistics is the smallest for X?
\\W) Range
\\X) Mean
\\Y) Median
\\Z) Mode
\\ANSWER: X) MEAN

\subsection*{BONUS}
\textbf{19) MATH} Short Answer Cailyn draws a circle, then draws another circle that has the same center but half the radius. She then repeats this process indefinitely, with each circle's radius equal to half the previous circle's radius. An ant is placed at random inside the largest circle, and begins to crawl in a straight line away from the center. What is the probability that the ant eventually crosses an odd number of circles?
\\ANSWER: 4/5 (ACCEPT: 0.8, 80%)

\subsection*{TOSS-UP}
\textbf{20) ENERGY} Short Answer The Van Voorhis Group can construct arbitrary molecular orbitals as linear combinations of predetermined linearly independent functions. What general term describes the set of these functions?
\\ANSWER: BASIS (ACCEPT: EIGENBASIS)

\subsection*{BONUS}
\textbf{20) ENERGY} Multiple Choice The Griffin Group developed a software package to extract information from electron paramagnetic resonance spectroscopy. Which of the following generates the signal of an EPR spectra?
\\W) Alignment of electron spin with an external electric field
\\X) Alignment of electron spin with a electric field generated by the nucleus
\\Y) Alignment of electron spin with an external magnetic field
\\Z) Alignment of electron spin with a magnetic field generated by the nucleus
\\ANSWER: Y) ALIGNMENT OF ELECTRON SPIN WITH AN EXTERNAL MAGNETIC FIELD

\subsection*{TOSS-UP}
\textbf{21) EARTH AND SPACE} Short Answer Order the following three celestial objects from least to greatest albedo: 1) Venus; 2) Jupiter; 3) Ceres.
\\ANSWER: 3, 2, 1

\subsection*{VISUAL BONUS}
\textbf{21) EARTH AND SPACE} Short Answer Your next bonus is visual. Answer the following two questions about the Earth-orbiting satellite whose ground track is shown in the image. 1) To the nearest hour, what is its period? 2) To the nearest 10 degrees, above what latitude does it reach apogee? Specify north or south.
\\ANSWER: 1) 12; 2) 60, NORTH (ACCEPT: 2) 60°N)

\subsection*{TOSS-UP}
\textbf{22) PHYSICS} Multiple Choice Stanley has a small bar magnet. He tries to create isolated north and south poles by slicing the bar magnet in half. Which of the following best describes the direction of the force between the two resulting pieces, and its power dependence on the distance r between the pieces?
\\W) Attractive, proportional to $1/r^{2}$
\\X) Attractive, proportional to $1/r^{4}$
\\Y) Repulsive, proportional to $1/r^{2}$
\\Z) Repulsive, proportional to $1/r^{4}$
\\ANSWER: X) ATTRACTIVE, PROPORTIONAL TO $1/r^{4}$ (READ: ONE OVER R TO THE FOURTH POWER)

\subsection*{BONUS}
\textbf{22) PHYSICS} Short Answer An infinitely long wire carries a current that is linearly increasing in time. Alborz has a rectangular loop, which he orients such that two of the sides are parallel to the wire. The closer of these two sides is 3 meters away from the wire, while the further side is 12 meters from the wire. If the rectangular loop is suddenly stretched such that the far side is now 24 meters from the wire, by what factor did the induced emf through the wire change?
\\ANSWER: 1.5 (ACCEPT: $3/2$)

\subsection*{TOSS-UP}
\textbf{23) BIOLOGY} Multiple Choice Dr. Damon is prescribing MAO (read: maow) inhibitors. Which of the following diseases is he likely treating?
\\W) ALS
\\X) Addisons
\\Y) Parkinson's
\\Z) ADHD
\\ANSWER: Y) PARKINSON'S

\subsection*{VISUAL BONUS}
\textbf{23) BIOLOGY} Short Answer Your next bonus is visual. A cross is made between an Hfr trp+ leu+ arg+ (read: HFR trip plus lew plus arj plus) bacteria and a F- trp-leu- arg- (read: F minus trip minus lew minus arj minus) bacteria and allowed to recombine. Given that the genes are in order and that arg (read: arj) enters last, order the following three resulting genotypes for the resulting arg+ F- (read: arj plus F minus) bacteria from least to most likely. 1) trp-leu- arg+ (read: trip minus lew minus arj plus); 2) trp+ leu- arg+ (read: trip plus lew minus arj plus); 3) trp+ leu+ arg+ (read: trip plus lew plus arj plus) recombination.
\\ANSWER: 2, 1, 3

\subsection*{TOSS-UP}
\textbf{24) CHEMISTRY} Short Answer What rule in thermodynamics is used to predict the absence of a quadruple point in a closed system containing only one pure component?
\\ANSWER: GIBBS PHASE RULE

\subsection*{BONUS}
\textbf{24) CHEMISTRY} Short Answer Identify all of the following three conditions that can be used to remove an acetal protecting group: 1) Acidic; 2) Basic; 3) Reductive.
\\ANSWER: 1 ONLY

\subsection*{TOSS-UP}
\textbf{25) PHYSICS} Short Answer By Noether's theorem, the first law of thermodynamics results from the invariance in translations of what quantity?
\\ANSWER: TIME (ACCEPT: TEMPORAL)

\subsection*{BONUS}
\textbf{25) PHYSICS} Short Answer According to Stoke's law, the drag force on a small sphere in a viscous fluid is linearly proportional to the product of the fluid's dynamic viscosity, the sphere's radius, and the sphere's velocity. What are the units of dynamic viscosity in SI base units?
\\ANSWER: KILOGRAM PER METER PER SECOND (ACCEPT: KILOGRAM INVERSE METER INVERSE SECOND)

\end{document}